\documentclass[10pt,a4paper]{amsart}
\usepackage{amsmath,amssymb,amsthm,graphicx}
\usepackage{verbatim}
\newenvironment{reference}{\comment}{\endcomment}
%\newenvironment{reference}{}{}
\newenvironment{slogan}{\comment}{\endcomment}
\newenvironment{history}{\comment}{\endcomment}

% For commutative diagrams we use Xy-pic
\usepackage[all]{xy}

% We use 2cell for 2-commutative diagrams.
\xyoption{2cell}
% Theorem environments.
%
\theoremstyle{plain}
\newtheorem{theorem}[subsection]{Theorem}
\newtheorem{proposition}[subsection]{Proposition}
\newtheorem{lemma}[subsection]{Lemma}

\theoremstyle{definition}
\newtheorem{definition}[subsection]{Definition}
\newtheorem{example}[subsection]{Example}
\newtheorem{exercise}[subsection]{Exercise}
\newtheorem{situation}[subsection]{Situation}

\theoremstyle{remark}
\newtheorem{remark}[subsection]{Remark}
\newtheorem{remarks}[subsection]{Remarks}

\numberwithin{equation}{subsection}

% Macros
%
\def\lim{\mathop{\mathrm{lim}}\nolimits}
\def\colim{\mathop{\mathrm{colim}}\nolimits}
\def\Spec{\mathop{\mathrm{Spec}}}
\def\Hom{\mathop{\mathrm{Hom}}\nolimits}
\def\Ext{\mathop{\mathrm{Ext}}\nolimits}
\def\SheafHom{\mathop{\mathcal{H}\!\mathit{om}}\nolimits}
\def\SheafExt{\mathop{\mathcal{E}\!\mathit{xt}}\nolimits}
\def\Sch{\mathit{Sch}}
\def\Mor{\mathop{\mathrm{Mor}}\nolimits}
\def\Ob{\mathop{\mathrm{Ob}}\nolimits}
\def\Sh{\mathop{\mathit{Sh}}\nolimits}
\def\NL{\mathop{N\!L}\nolimits}
\def\CH{\mathop{\mathrm{CH}}\nolimits}
\def\proetale{{pro\text{-}\acute{e}tale}}
\def\etale{{\acute{e}tale}}
\def\QCoh{\mathit{QCoh}}
\def\Ker{\mathop{\mathrm{Ker}}}
\def\Im{\mathop{\mathrm{Im}}}
\def\Coker{\mathop{\mathrm{Coker}}}
\def\Coim{\mathop{\mathrm{Coim}}}

% Boxtimes
%
\DeclareMathSymbol{\boxtimes}{\mathbin}{AMSa}{"02}

%
% Macros for moduli stacks/spaces
%
\def\QCohstack{\mathcal{QC}\!\mathit{oh}}
\def\Cohstack{\mathcal{C}\!\mathit{oh}}
\def\Spacesstack{\mathcal{S}\!\mathit{paces}}
\def\Quotfunctor{\mathrm{Quot}}
\def\Hilbfunctor{\mathrm{Hilb}}
\def\Curvesstack{\mathcal{C}\!\mathit{urves}}
\def\Polarizedstack{\mathcal{P}\!\mathit{olarized}}
\def\Complexesstack{\mathcal{C}\!\mathit{omplexes}}
% \Pic is the operator that assigns to X its picard group, usage \Pic(X)
% \Picardstack_{X/B} denotes the Picard stack of X over B
% \Picardfunctor_{X/B} denotes the Picard functor of X over B
\def\Pic{\mathop{\mathrm{Pic}}\nolimits}
\def\Picardstack{\mathcal{P}\!\mathit{ic}}
\def\Picardfunctor{\mathrm{Pic}}
\def\Deformationcategory{\mathcal{D}\!\mathit{ef}}
\usepackage{cleveref}
\crefname{originallemma}{original lemma}{original lemmas}
\Crefname{originallemma}{Original lemma}{Original lemmas}
\makeatletter
\newlabel{lemma-ideals-valuation-ring}{{00IH}{source}{}{sourcepointer-1}{}}
\newlabel{lemma-ideals-valuation-ring@cref}{{[originallemma][0][]00IH}{[1][0][]source}{}{}{}}
\makeatother
\title{A valuation ring is Noetherian exactly when it is a discrete valuation ring or a field}
\author{The Stacks Project Authors}
\date{stacks-project:algebra.tex:local-snapshot-20261012; GNU FDL 1.2 or later}
\begin{document}
\maketitle
\noindent\textit{Editorial note.} All delivered mathematical text is the author's own, copied from the pinned Stacks Project chapter file algebra.tex; no statement or proof step was rewritten, reordered or removed.\par\medskip
\noindent\textit{Editorial note.} Preamble provenance, as required for this source: the frozen snapshot stages the chapter's own preamble file beside it (preamble.tex, 217 lines, 8116 bytes, the same shared Stacks-wide preamble the published chow and etale-cohomology cards use); the first line of the pinned chapter file is \textbackslash{}input\{preamble\}, so no preamble text lives inside the chapter itself. The card takes that file's lines 8-18 (the verbatim package and the three \textbackslash{}comment-based environments reference/slogan/history at lines 9-12, section 13.2, together with the Xy-pic setup) and its lines 154 to the runtime end of the file (the native theorem declarations and mathematical macros). Not one character was changed, no package was added, no macro was defined, and the excerpt serves only as preamble, never as a span coordinate.\par\medskip
\noindent\textit{Editorial note.} The same-chapter pointers are located at their \textbackslash{}label lines inside the pinned chapter file; the cross-chapter pointers are located in the project's own tag list staged in the snapshot (sources/stacks-chow/tags.txt), at the verbatim lines carrying their canonical tags.\par\medskip
\noindent\textit{Editorial note.} The counted proof refers to the correspondence between ideals and ideals of the value group by label; it is recorded as a navigation stub pointed at its \textbackslash{}label line, and the argument consuming it stays inside the counted span.\par\medskip
\noindent\textit{Editorial note.} The counted spans contain no \textbackslash{}cite, so the card carries no bibliography.\par\medskip
\noindent\textit{Allowed standard prerequisites (editorial).}\begin{itemize}
\item Ideals of a valuation ring correspond to ideals of the value group. Used in both halves of the counted proof (lines 12323 and 12330). Translates the ideal-theoretic chain conditions into statements about the value group.
\item The value group of a discrete valuation ring is the integers. Used in the first half of the counted proof (lines 12320-12325). Produces explicit generators of the ideals and hence a principal ideal domain.
\end{itemize}\medskip
\ifdefined\hypertarget\hypertarget{sourcepointer-1}{}\fi
\noindent\textit{Source reference.} Original lemma 00IH: sources/stacks-algebra/pinned-original-source.tex, source lines 12302--12302. Navigation pointer for \textbackslash{}ref\{lemma-ideals-valuation-ring\} used inside the counted span (line 12323); the argument that consumes it is retained in the span.\par\medskip
\begin{lemma}
\label{lemma-valuation-ring-Noetherian-discrete}
A valuation ring is Noetherian if and only if it is
a discrete valuation ring or a field.
\end{lemma}

\begin{proof}
Suppose $A$ is a discrete valuation ring
with valuation $v : A \setminus \{0\} \to \mathbf{Z}$
normalized so that $\Im(v) = \mathbf{Z}_{\geq 0}$.
By Lemma \ref{lemma-ideals-valuation-ring} the ideals of $A$ are the subsets
$I_n = \{0\} \cup v^{-1}(\mathbf{Z}_{\geq n})$. It is clear
that any element $x \in A$ with $v(x) = n$ generates $I_n$.
Hence $A$ is a PID so certainly Noetherian.

\medskip\noindent
Suppose $A$ is a Noetherian valuation ring with value group $\Gamma$.
By Lemma \ref{lemma-ideals-valuation-ring} we see the ascending chain
condition holds for ideals in $\Gamma$.
We may assume $A$ is not a field, i.e., there is a $\gamma \in \Gamma$
with $\gamma > 0$. Applying the ascending chain condition to the subsets
$\gamma + \Gamma_{\geq 0}$ with $\gamma > 0$ we see
there exists a smallest element $\gamma_0$ which is bigger than $0$.
Let $\gamma \in \Gamma$ be an element $\gamma > 0$. Consider the sequence
of elements $\gamma$, $\gamma - \gamma_0$, $\gamma - 2\gamma_0$,
etc. By the ascending chain condition these cannot all be $> 0$.
Let $\gamma - n \gamma_0$ be the last one $\geq 0$. By minimality
of $\gamma_0$ we see that $0 = \gamma - n \gamma_0$. Hence $\Gamma$
is a cyclic group as desired.
\end{proof}

\end{document}
